% !TEX root = ../main.tex
\begin{table}[H]
    \centering
    \caption{Modele rozważane w projekcie}
    \label{tab:kandydaci_modeli}

    \small
    \setlength{\tabcolsep}{7pt}
    \renewcommand{\arraystretch}{1.35}
    \begin{tabular}{|p{3.8cm}|p{2.3cm}|p{2.3cm}|>{\raggedright\arraybackslash}p{5.3cm}|}
        \hline
        \rowcolor{gray!20}
        \textbf{Model} & \textbf{Twórca} & \textbf{Parametry} & \textbf{Decyzja i powód} \\ \hline
        \textbf{gemma4:12b} &
        \textbf{Google} &
        \textbf{12 mld} &
        \textbf{Wybrany.} Najnowsza generacja rodziny Gemma, dobra jakość na zadaniach matematycznych, mieści się w limicie pamięci. \\ \hline
        \textbf{qwen3.5:9b} &
        \textbf{Alibaba} &
        \textbf{9 mld} &
        \textbf{Wybrany.} Najnowsza generacja rodziny Qwen, dobra jakość na zadaniach matematycznych, inna rodzina niż Gemma. \\ \hline
        llama3.1:8b &
        Meta &
        8 mld &
        Odrzucony. Wyraźnie słabsza jakość odpowiedzi. \\ \hline
        qwen2.5:7b &
        Alibaba &
        7 mld &
        Odrzucony. Jakość przeciętna oraz starsza generacja względem Qwen3.5. \\ \hline
        gemma2:9b &
        Google &
        9 mld &
        Odrzucony. Jakość przeciętna oraz starsza generacja względem Gemma 4. \\ \hline
        gemma2:2b &
        Google &
        2 mld &
        Odrzucony. Częstsze błędy i trudności z zachowaniem wymaganego formatu odpowiedzi. \\ \hline
        phi4-mini &
        Microsoft &
        3,8 mld &
        Odrzucony. Częstsze błędy i trudności z zachowaniem wymaganego formatu odpowiedzi. \\ \hline
        deepseek-r1:14b &
        DeepSeek-AI &
        14 mld &
        Odrzucony. Bardzo dobre wyniki w matematyce, zbliżone do wybranych modeli, ale wyższe zużycie pamięci. \\ \hline
    \end{tabular}
\end{table}
